\chapter{Inventar de endpoint-uri API}

\section{Observații generale}

Această anexă sintetizează endpoint-urile principale ale aplicației. Nu urmăresc să refac documentația Swagger în întregime, ci să ofer o imagine clară asupra suprafeței API-ului și a modului în care aceasta este împărțită pe domenii funcționale.

\section{Endpoint-uri pentru cont și autentificare}

\begin{longtable}{>{\RaggedRight\arraybackslash}p{1.7cm}>{\RaggedRight\arraybackslash}p{6.1cm}>{\RaggedRight\arraybackslash}p{5.8cm}}
  \caption{Endpoint-uri principale din zona de autentificare} \\
  \toprule
  \textbf{Metodă} & \textbf{Rută} & \textbf{Rol} \\
  \midrule
  \endfirsthead
  \toprule
  \textbf{Metodă} & \textbf{Rută} & \textbf{Rol} \\
  \midrule
  \endhead
  POST & \path{/api/account/register} & creare cont nou \\
  POST & \path{/api/account/login} & autentificare clasică \\
  POST & \path{/api/account/oauth2/github} & autentificare prin GitHub \\
  POST & \path{/api/account/oauth2/google} & autentificare prin Google \\
  POST & \path{/api/account/two-factor/verify-login} & finalizarea login-ului cu 2FA \\
  POST & \path{/api/account/oauth2/github/link} & legare cont GitHub la utilizatorul curent \\
  POST & \path{/api/account/oauth2/google/link} & legare cont Google la utilizatorul curent \\
  POST & \path{/api/account/forgot-password} & inițiere resetare parolă \\
  POST & \path{/api/account/reset-password} & setarea noii parole prin token \\
  POST & \path{/api/account/set-password} & adăugarea unei parole pentru cont OAuth \\
  POST & \path{/api/account/change-password} & schimbarea parolei \\
  POST & \path{/api/account/two-factor/enable/request} & trimiterea codului pentru activare 2FA \\
  POST & \path{/api/account/two-factor/enable/confirm} & confirmarea activării 2FA \\
  POST & \path{/api/account/two-factor/disable/request} & trimiterea codului pentru dezactivare 2FA \\
  POST & \path{/api/account/two-factor/disable/confirm} & confirmarea dezactivării 2FA \\
  POST & \path{/api/account/refresh} & reînnoirea token-ului de acces \\
  POST & \path{/api/account/logout} & încheierea sesiunii \\
  \bottomrule
\end{longtable}

\section{Endpoint-uri pentru proiecte}

\begin{longtable}{>{\RaggedRight\arraybackslash}p{1.7cm}>{\RaggedRight\arraybackslash}p{6.1cm}>{\RaggedRight\arraybackslash}p{5.8cm}}
  \caption{Endpoint-uri principale pentru proiecte} \\
  \toprule
  \textbf{Metodă} & \textbf{Rută} & \textbf{Rol} \\
  \midrule
  \endfirsthead
  \toprule
  \textbf{Metodă} & \textbf{Rută} & \textbf{Rol} \\
  \midrule
  \endhead
  GET & \path{/api/projects} & lista proiectelor utilizatorului curent \\
  GET & \path{/api/projects/{id}} & proiect după identificator \\
  GET & \path{/api/projects/by-slug/{slug}} & proiect după slug \\
  POST & \path{/api/projects} & creare proiect \\
  PUT & \path{/api/projects/{id}} & actualizare metadate proiect \\
  DELETE & \path{/api/projects/{id}} & ștergere proiect \\
  POST & \path{/api/projects/{id}/assets/images} & upload imagine pentru canvas \\
  GET & \path{/api/projects/{id}/design} & încărcare design salvat \\
  PUT & \path{/api/projects/{id}/design} & salvare design \\
  POST & \path{/api/projects/{id}/flush} & salvare combinată design + thumbnail \\
  PUT & \path{/api/projects/{id}/thumbnail} & actualizare thumbnail \\
  GET & \path{/api/projects/user/{userId}} & proiecte publice ale unui utilizator \\
  POST & \path{/api/projects/{id}/star} & bookmark / star \\
  DELETE & \path{/api/projects/{id}/star} & eliminare bookmark \\
  POST & \path{/api/projects/{id}/view} & înregistrare vizualizare \\
  POST & \path{/api/projects/{id}/like} & apreciere proiect \\
  DELETE & \path{/api/projects/{id}/like} & eliminare apreciere \\
  POST & \path{/api/projects/{id}/fork} & creare fork după proiect public \\
  \bottomrule
\end{longtable}

\section{Endpoint-uri pentru utilizatori și profiluri}

\begin{longtable}{>{\RaggedRight\arraybackslash}p{1.7cm}>{\RaggedRight\arraybackslash}p{6.1cm}>{\RaggedRight\arraybackslash}p{5.8cm}}
  \caption{Endpoint-uri principale pentru utilizatori} \\
  \toprule
  \textbf{Metodă} & \textbf{Rută} & \textbf{Rol} \\
  \midrule
  \endfirsthead
  \toprule
  \textbf{Metodă} & \textbf{Rută} & \textbf{Rol} \\
  \midrule
  \endhead
  GET & \path{/api/users/me} & profilul utilizatorului curent \\
  PUT & \path{/api/users/me} & actualizare profil propriu \\
  POST & \path{/api/users/me/profile-image} & upload imagine de profil \\
  DELETE & \path{/api/users/me} & ștergere cont propriu \\
  DELETE & \path{/api/users/me/linked-accounts/{provider}} & deconectare provider OAuth \\
  GET & \path{/api/users/search?q=...} & căutare de utilizatori \\
  GET & \path{/api/users/{username}} & profil public după username \\
  POST & \path{/api/users/{username}/follow} & urmărire utilizator \\
  DELETE & \path{/api/users/{username}/follow} & oprire urmărire \\
  GET & \path{/api/users/{username}/followers} & lista urmăritorilor \\
  GET & \path{/api/users/{username}/following} & lista utilizatorilor urmăriți \\
  GET & \path{/api/users/me/stars} & proiectele salvate de utilizatorul curent \\
  \bottomrule
\end{longtable}

\section{Endpoint-uri pentru explore și recomandări}

\begin{longtable}{>{\RaggedRight\arraybackslash}p{1.7cm}>{\RaggedRight\arraybackslash}p{6.1cm}>{\RaggedRight\arraybackslash}p{5.8cm}}
  \caption{Endpoint-uri din zona explore} \\
  \toprule
  \textbf{Metodă} & \textbf{Rută} & \textbf{Rol} \\
  \midrule
  \endfirsthead
  \toprule
  \textbf{Metodă} & \textbf{Rută} & \textbf{Rol} \\
  \midrule
  \endhead
  GET & \path{/api/explore/trending} & proiecte populare \\
  GET & \path{/api/explore/recommended} & recomandări de proiecte, eventual personalizate \\
  GET & \path{/api/explore/people} & sugestii de utilizatori \\
  \bottomrule
\end{longtable}

\section{Endpoint-uri pentru AI și conversie}

\begin{longtable}{>{\RaggedRight\arraybackslash}p{1.7cm}>{\RaggedRight\arraybackslash}p{6.1cm}>{\RaggedRight\arraybackslash}p{5.8cm}}
  \caption{Endpoint-uri pentru AI și converter} \\
  \toprule
  \textbf{Metodă} & \textbf{Rută} & \textbf{Rol} \\
  \midrule
  \endfirsthead
  \toprule
  \textbf{Metodă} & \textbf{Rută} & \textbf{Rol} \\
  \midrule
  \endhead
  POST & \path{/api/ai/design} & generare design într-o singură cerere \\
  POST & \path{/api/ai/design/stream} & generare design în regim SSE \\
  POST & \path{/api/ai/design/pipeline} & rulare pipeline AI în trei faze \\
  POST & \path{/api/ai/design/pipeline/stream} & rulare pipeline AI cu streaming \\
  POST & \path{/api/converter/generate} & export pentru o pagină sau set responsive \\
  POST & \path{/api/converter/validate} & validare structură IR \\
  POST & \path{/api/converter/generate-files} & generare fișiere pentru export multi-page \\
  \bottomrule
\end{longtable}

\noindent\textbf{Observație finală.}

Din inventarul de mai sus se vede că API-ul nu este limitat la o singură zonă funcțională. El trebuie să deservească atât fluxuri de produs clasice, cât și funcționalități tehnice speciale, cum sunt streaming-ul AI și exportul multi-framework. Tocmai această diversitate justifică structura pe straturi din backend și separarea convertorului în proiect propriu.
